\chapter{电子不是绕核飞行的小行星}

\begin{kaichang}
很多城市的老街上还挂着一种老式路灯。天一黑，它们先发出暗红色的光，几分钟以后慢慢变成明亮的橙黄色，把整条街照得像泡在橘子汽水里。这种灯叫\textbf{钠灯}，里面装的是一点点钠金属蒸气和惰性气体。

这里有个很值得想一想的问题：霓虹灯可以做成红的、蓝的、绿的、紫的，五光十色；可钠灯永远是那一种橙黄色，谁也换不掉。是厂家偷懒吗？不是。橙黄色不是油漆，不是玻璃颜色，而是从钠原子内部“长”出来的颜色——钠只会发这个颜色，也只会发这个颜色，就像一个人只会唱某几个固定的音。

再看看节日的烟花：红色、绿色、金色、蓝色，又是谁“染”上去的？先猜一个答案写下来。这一章结束时，你会知道颜色背后藏着一套严格的“楼层住宿规则”，而这套规则正是下一张大地图——元素周期表——的地基。
\end{kaichang}

\section{把光拆开看一看}

颜色是可以“拆开”的。阳光穿过三棱镜，会展开成一条连续的彩虹，从红到紫，什么颜色都有，一段不缺。牛顿当年就是这样证明白光是由各种颜色混合成的。你也可以在CD光盘背面、在肥皂泡上看到类似的连续色彩。

现在换个光源。把钠灯的光送进棱镜（物理实验室里用更精细的“分光镜”），你猜会看到什么？不是彩虹。是几条孤立的亮线：两条挨得很近的黄色亮线特别刺眼，其余地方几乎是黑的。霓虹灯里的氖气发红光，拆开看也是一组分离的亮线，不是连续彩带。每一种气体、每一种元素，拆开后的“线条图案”都不同，而且无论你在北京还是巴黎测量，无论这团气体在灯管里还是在恒星上，图案都一模一样。

这件事在19世纪就把科学家逼到了墙角：\textbf{每种元素都有一套固定不变的“光谱指纹”}。光是从原子内部出来的，指纹固定，说明原子内部一定存在某种固定的、不连续的结构。看不见的原子，居然通过光把自己的内部结构“寄”给了我们。

这套指纹今天就在你身边工作。老式日光灯管里，汞蒸气先发出紫外线，管壁上的荧光粉再把紫外线“翻译”成几种可见光混成白色；验钞机用紫外线激发纸币上的荧光标记，标记发光的颜色也是原子指纹；天文馆门口卖的“光谱眼镜”，对着路灯一照就能看到它属于钠灯还是LED灯——两种“白色”拆开后的图案完全不同。看颜色，是分辨物质最便宜也最快的一招。

还有一个更贴近生活的版本。往火焰里送进一点钠盐，火苗立刻变成明亮的黄色——这就是路灯那种黄。这个“焰色”现象你在家里就能安全地试一次。

\begin{shiyan}
\textbf{观察食盐的焰色}（建议有成年人在场）

\textbf{材料}：食盐一小勺、清水小半杯、一根棉签（或一把干净的不锈钢勺）、一支蜡烛或燃气灶的小火。

\textbf{步骤}：把食盐溶进水里，搅成浓盐水。点燃蜡烛，先看一眼火焰本来的颜色（淡蓝色为主，烛芯处发黄）。用棉签蘸饱盐水，把湿棉签头伸到火焰上缘，别碰到烛芯，观察火焰颜色的变化。多蘸几次盐水，重复观察。

\textbf{观察点}：①蘸盐水的棉签让火焰变成什么颜色？②是纯黄色还是杂色？③把棉签拿走，颜色是否立刻消失？

\textbf{安全提示}：棉签蘸湿再靠近火焰，烧着了就浸入水杯中；不要往火焰里撒干盐；结束后熄灭蜡烛、收拾好火柴。只用食盐，不要尝试其他粉末。

你看到的亮黄色，和钠灯、和烟花里的金色，是同一种原子发出来的同一组光谱线。注意：食盐是氯化钠（\chem{NaCl}），这抹黄色属于钠，不属于氯——换成氯化钾，火焰会变成淡紫色。火焰就像一个临时的“原子广播站”，把钠的身份用颜色播报了出来。
\end{shiyan}

\section{行星模型撞上两堵墙}

第7章讲过，卢瑟福用金箔实验发现：原子的几乎全部质量和全部正电荷都缩在一个极小的核里，电子在核外。一个自然的想象马上冒出来：电子像行星绕太阳一样，绕着核转圈。这个“行星模型”好懂、好画，直到今天还印在无数课本的封面上。

可惜，它错得很干脆。两堵墙挡在前面。

第一堵墙来自电磁学。19世纪的物理已经算明白：带电粒子只要做加速运动（转圈就是不停改变方向的加速运动），就会向外辐射电磁波、不断损失能量。如果电子真像行星一样绕核转，它应该一边转一边发光、一边丢失能量，轨道越转越小，最后螺旋着一头栽进原子核。算出来的崩塌时间大约是 $10^{-10}$ 秒——也就是说，行星模型里的原子连一亿分之一秒都活不过去。可你手中的这本书、这本书里的每一个原子，都稳稳当当地存在着。这是行星模型的死罪：它预言的原子根本站不住。

第二堵墙来自上一节的光谱。行星可以停在任何半径的轨道上，能量可以连续变化，所以辐射出来的光应该是连续的一片——可实际测到的偏偏是分离的亮线。电子如果随便什么能量都能取，光谱就不该有“指纹”。

矛盾摆在桌面上：原子稳定存在，光谱线固定分立。这两个事实合在一起，等于大自然在宣布：电子的能量不能连续变化，\textbf{只能取某些特定的值}。就像楼梯——你可以站在这一级或那一级，但没法站在两级之间的半空中。

\section{电子住的是“能量楼层”，不是轨道}

1913年，玻尔给氢原子立下一条“蛮横”的新规矩：电子只能待在一组特定的能级上（能级就是“允许存在的能量值”），待在里面时不辐射能量；只有当电子从一个能级跳到另一个能级时，才会吸收或吐出一个光子，光子的能量恰好等于两级之差。

这一条立刻解释了一切。为什么光谱是分离的线？因为能级差是固定的，吐出的光子能量就那么几种，颜色自然只有那么几种。为什么每种元素指纹不同？因为不同元素的原子核带的正电荷不同，能级的排布就不同，像每栋楼的台阶高度都不一样。为什么原子不会塌？因为存在一个最低的能级（叫\textbf{基态}，就是“能量最低的住宿层”），电子待在那里时没有更低的层可跳，无处可塌。

\begin{tikzpicture}[scale=0.95]
\draw[thick] (0,0) -- (3,0) node[right] {$n=1$（基态，能量最低）};
\draw[thick] (0,1.3) -- (3,1.3) node[right] {$n=2$};
\draw[thick] (0,2.3) -- (3,2.3) node[right] {$n=3$};
\draw[thick] (0,3.1) -- (3,3.1) node[right] {$n=4$};
\node[right] at (3,3.9) {越往上能级越挤，再往上就“自由”了};
\draw[->,thick] (2.2,2.3) -- (2.2,0.08);
\node[right] at (2.25,1.2) {跳下来：吐出一个光子};
\draw[->,thick] (0.8,0.08) -- (0.8,1.28);
\node[right] at (0.85,0.7) {跳上去：先吸收能量};
\end{tikzpicture}

上面这张图是人为的表示法：横线代表允许的能量值，不代表电子在空间里真的站在某条线上。能级是\textbf{能量的刻度}，不是地理位置的楼层。

电子在高层和低层之间跳上跳下，听起来像粒子的杂技。但真正的麻烦还在后面：电子在某一能级上的时候，到底“在”哪儿？

\section{电子云：一张概率地图，不是运动残影}

1920年代，薛定谔、海森堡等人建立的量子力学给出了让所有人都不舒服、但至今没被任何实验推翻过的答案：\textbf{电子在两次测量之间，没有确定的位置，也没有轨迹}。

请把这句话读慢一点。它不是“电子跑得太快，我们看不清它在哪”，也不是“仪器不够精密，将来就能测准”。它的意思是：在测量之前，“电子此刻在哪个点”这个问题本身就没有答案。量子力学能算的，是你在核外某个位置找到电子的\textbf{概率}——离核近的地方概率高还是低、往哪个方向概率高，这些可以算得极准，准到和实验小数点后好多位都对得上。

“轨道”（orbital）这个词就是这样被保留下来的，但它的意思已经彻底换了：一个轨道不是一条路线，而是\textbf{一个概率分布的形状}——一张“电子容易在哪儿被找到”的三维地图。把概率高低用疏密不同的点画出来，就得到常说的\textbf{电子云}：

\begin{itemize}[itemsep=2pt]
\item \textbf{s 轨道}：球形。以核为中心，某个距离附近最容易找到电子，方向上没有偏好。
\item \textbf{p 轨道}：哑铃形，两瓣，分别伸在核的两侧；三个 p 轨道沿 $x$、$y$、$z$ 三个互相垂直的方向各摆一副。
\item \textbf{d、f 轨道}：形状更复杂，花瓣更多。中学阶段知道它们存在就够了。
\end{itemize}

这些形状都是概率地图的轮廓，不是硬壳，不是气球，不是电子“占着”的一块区域。画成图时我们常常会圈出一条边界，那只是“电子有九成概率出现在这条线以内”的参考线，就像天气预报画的降雨范围，不是一堵墙。

那“概率”到底该怎么理解？想象这样一个实验：你有一台假想的超级仪器，能一次次地“抓拍”氢原子基态电子的位置。每按一次快门，照片上都只有一个孤零零的点——电子从来不是一小片云，每次出现都是一个点。拍一万次，把一万个点叠在同一张图上：点多的地方连成浓密的影，点少的地方稀稀落落，整体呈现一个以核为中心的球形图案。电子云就是这张叠加图的样子。\textbf{云是无数次定位的统计结果，不是某一次观测中电子摊开成了一朵云，更不是电子高速飞奔拖出来的模糊轨迹。}

\begin{wuqu}
“电子云是电子绕核飞得太快，像转动的扇叶连成圆盘那样，拖出来的残影。”——这个误解流传极广，因为它看起来好懂。但它从根上错了。

反例和反驳有三条。第一，实验里每次探测电子，它都出现在一个确定的点，从来没有人拍到过一小段“拉长的轨迹”；如果真有轨迹，能量足够高时理应能看出扫动的痕迹，可是没有。第二，运动残影解释不了稳定性：一个真有轨道、真在加速飞奔的电子，照电磁学规律就该辐射能量、螺旋坠核（本章前面算过，约 $10^{-10}$ 秒），原子早就塌了；电子“不塌”恰恰因为它根本没有那种绕圈运动。第三，基态氢原子的电子云是球对称、不随时间变化的静止图案——一朵“飞奔留下的残影”凭什么永远均匀、永远不变？

正确的图景是：电子没有轨迹可拖，云图描绘的是“在不同位置找到它的概率”，概率本身就长那个形状。承认这一点确实别扭，但一百年来所有实验都在为这幅别扭的图景投票。
\end{wuqu}

顺带说一句行星模型真正的死因，其实比“会坠核”更深：在量子力学里，一个被限制在原子这么小范围里的电子，根本不可能有确定的速度和轨道——这不是技术限制，而是自然规律（海森堡不确定性原理）。电子越小越“宅”在原子里，它的“运动状态”就越模糊。想让电子乖乖走一条确定的圆轨道，等于要求它同时拥有确定的位置和确定的速度，大自然不答应。

\section{楼层旅馆的住宿规则}

多电子原子有几十个电子要安排，大自然有一套严格的“旅馆住宿规则”。这套规则不需要背，理解三条就够了。

\textbf{第一条：楼层和户型。}每个电子的能量主要由主量子数 $n$ 决定，$n=1,2,3,\cdots$ 就是第几层，数字越大能量越高、离核的平均距离越远。同一层里还有不同“户型”：s、p、d、f，对应不同形状的概率地图。第1层只有 s 户型（1s），第2层有 s 和 p（2s、2p），第3层再添 d，第4层再添 f。

\textbf{第二条：每间房最多住两人，且必须“背靠背”。}电子除了位置和能量，还有一个像小陀螺一样的内禀属性叫\textbf{自旋}（它并不是真的在自转，只是表现得像有两种朝向）。\textbf{泡利不相容原理}说：同一个轨道里最多容纳两个电子，而且这两个电子的自旋必须相反。所以 1s 最多 2 个，三个 2p 轨道合起来最多 6 个，第2层总共 8 个——楼层容量 $2,8,18,32$ 就是这么来的。

\textbf{第三条：先住便宜的，同房型先一人一间。}电子总是先占能量最低的轨道，这叫\textbf{构造原理}：1s 住满才轮到 2s，2s 满了才轮到 2p。同一户型的几个房间（比如三个 2p）能量一样，电子会先一人一间、且自旋同向地分开住，住不下才开始两两合住——这叫\textbf{洪特规则}。直觉类比：长椅上的陌生人总是先一人坐一条，坐满了才愿意挨着坐。背后的原因是同种电荷互相排斥，分开住能量更低。

现在轮到符号登场。用“楼层+户型+人数”的写法，电子的住宿表一目了然：氢是 \chem{1s^1}，氦是 \chem{1s^2}（第一层住满），碳是 \chem{1s^2 2s^2 2p^2}，钠是 \chem{1s^2 2s^2 2p^6 3s^1}。看钠：前两层像氖一样住得满满当当，只有第3层孤零零 1 个电子——这颗“独苗电子”几乎决定了钠的全部脾气。至于这些脾气怎样让元素们排成一张周期表，第10章会带你把整张地图铺开。

\textbf{规矩也有例外，而且不许掩盖。}按构造原理，铬应该是 \chem{3d^4 4s^2}，实测却是 \chem{3d^5 4s^1}；铜应该是 \chem{3d^9 4s^2}，实测是 \chem{3d^{10} 4s^1}。半满（5个）和全满（10个）的 d 户型格外“划算”，一个 4s 电子就搬了家。这不是规则失效，而是提醒我们：住宿规则本身是近似，真实的能量账要细到电子之间每一份排斥。科学里“规则+诚实的例外”远比“完美但虚假的口诀”可靠。

\begin{qiaoliang}
光子的能量由颜色决定：$E = hc/\lambda$，其中 $\lambda$ 是波长，$h$ 和 $c$ 是常数。波长越短（越偏蓝紫），每个光子的能量越高。

来算一笔真账。钠灯那两条黄线的波长约为 \ud{589}{nm}（$5.89\times10^{-7}$ 米），代入公式，每个光子携带的能量约 $3.4\times10^{-19}$ 焦耳——小得可怜。但看看数量级的另一端：一盏 \ud{100}{W} 的钠灯，假设电能全变成这种黄光，每秒要发出
\[100 \div (3.4\times10^{-19}) \approx 3\times10^{20}\ \text{个光子}。\]
每秒三十万亿亿个光子，每个都是某个钠原子里的电子从 3p 轨道跳回 3s 轨道时吐出来的。原子很小，但能级差换算成光，足以照亮整条街。

反着用同一个公式更有意义：\textbf{测出光的颜色，就能反推出原子内部的能级差}。这就是为什么光谱能成为探测原子内部的望远镜——不用切开原子，光自己把能级的间距报了出来。
\end{qiaoliang}

\section{先在太阳上发现的元素}

\begin{zhengju}
这套“指纹”想法不是谁拍脑袋想出来的，它有一段环环相扣的证据史。

1814年，德国镜片工匠夫琅禾费在测试棱镜时，让阳光穿过一条细缝再进棱镜，意外看到彩虹里压着几百条细细的\textbf{暗线}——某些特定颜色的光在到达地球的路上“失踪”了。他逐个记下暗线的位置，却解释不了。这个谜挂了四十多年。

1859年，海德堡的化学家本生和物理学家基尔霍夫把谜底揭开了。本生改装的煤气灯（就是今天实验室的本生灯）火焰温度高、本身几乎无色，把各种盐送进火焰，每种元素都烧出一套自己的亮线，图案绝不混淆。基尔霍夫接着做了关键一步：让强光先穿过钠蒸气再进棱镜，结果在原本连续的光谱上、恰好是钠亮线的位置，出现了暗线。结论：原子既能发射也能吸收自己那套特定颜色的光，太阳上的暗线，正是阳光穿过太阳外层气体时被那里的原子“截留”的指纹。光谱不仅能查地球上的元素，还能查九千多万公里外恒星的成分——当时有人刚断言“恒星的化学成分人类永远不可能知道”，话音未落就被推翻。

这套方法立刻变成了“元素探测器”。1860年，本生在矿泉水的残渣里发现两条从没见过的天蓝色亮线，宣布发现新元素，命名为铯（拉丁语“天蓝”）；第二年又发现了铷（“暗红”）。更大胆的一幕发生在1868年：法国天文学家让桑在日食时观测太阳边缘的日珥，发现一条黄色亮线，和钠的黄线位置对不上，和任何已知元素都对不上。英国天文学家洛克耶断言：这是太阳上存在而地球上尚未发现的一种新元素，取名氦（来自希腊语“太阳”）。不少同行觉得这太离谱——凭一条线就宣布新元素？二十七年后的1895年，拉姆齐在处理铀矿放出的气体时，光谱里出现了同一条线。氦，先在太阳上被发现，然后才在地球上找到。

剩下最后一个“为什么”：指纹为什么恰好是这些波长？1913年玻尔的能级模型算准了氢的光谱线，1920年代量子力学把多电子原子也装了进去。证据链的每一环——暗线、焰色、预言新元素、能级理论——都能被后来者检验和推翻，而它们全都站住了。这才是“科学家怎样知道”的标准姿势。
\end{zhengju}

\section{这幅图景的边界}

电子云模型极其成功，但它有几条明确的边界，知道边界才算真懂。

第一，玻尔的“能级轨道”混合模型只对氢原子精确有效，多电子原子必须靠量子力学近似计算——电子互相排斥，几万项的账只能交给计算机。我们嘴里说“楼层”，心里要记得那只是 $n$ 的近似含义：在多电子原子里，同一层的 s、p、d 能量并不相同，4s 甚至比 3d 更早被填入，铬和铜的例外就是这么来的。

第二，电子云图里的形状（球、哑铃）严格说是“一个电子”的概率地图。把它当成几十电子原子里每个电子各自的“房间”，是极好用的近似，但电子之间互相推挤，真实图景比单人房间复杂。

第三，概率描述不是“没办法的办法”。曾有人设想：电子其实有确定轨迹，只是我们不知道（所谓“隐变量”）。20世纪后半叶的贝尔不等式实验判决性地排除了这类解释中最自然的一大类——电子的“模糊”是大自然的本性，不是我们的无知。这条边界今天仍在被更精密的实验反复确认。

第四，别忘了电子云只管核外。原子核内部（质子、中子、放射性）是另一套规则，第8章已经讲过；而电子怎样在原子之间“谈判”、搭出化学键，要到第17章才开场。

\section{回到那盏橙黄色的路灯}

现在可以回答开场的问题了。钠灯里，电流把钠原子的外层电子从 3s 轨道推上 3p 轨道；电子跳回来时，吐出的光子能量恰好对应 \ud{589}{nm}——恰好是人眼最敏感的黄色。这个颜色换不掉，因为它是钠原子内部两级台阶的高度差决定的；台阶不挪，颜色不改。想让灯变颜色，唯一的办法是换元素：氖气发红橙光，氩气发蓝紫光，氙气发白光——满街的霓虹，其实是一排原子的“能级展览”。

烟花同理：锶盐烧出大红，钡盐烧出草绿，铜盐烧出湛蓝，而那种耀眼的金色，多半还是老熟人钠。焰色实验里你亲手“点亮”的，和路灯、烟花、恒星光谱里的，是同一条物理规律的三个现场。

再深一层：每一种颜色对应一个确定的能级差，而能级差的整条阶梯，就是那套楼层住宿规则的直接后果。下一章你会看到，这套规则最惊人的产物不是颜色，而是一张表：为什么表格的宽度恰好是 2、8、8、18？为什么同一列的元素脾气像一个模子刻出来的？为什么门捷列夫敢在表里留空格、预言还没被发现的元素？第10章，我们把电子的住宿表摊开，看它怎样长成元素周期表那座造型古怪的城市。

\begin{tiaozhan}
玻尔模型里最漂亮的战果，是一个公式算出氢原子的全部能级：$E_n = -13.6/n^2$（单位是电子伏特 eV，$n=1,2,3,\cdots$）。负号表示电子被核“拴着”，能量比自由状态低。

试试用它预言一条光谱线。电子从 $n=3$ 跳到 $n=2$，吐出的光子能量是
\[\Delta E = 13.6\times\left(\frac{1}{2^2}-\frac{1}{3^2}\right) = 13.6\times\frac{5}{36} \approx 1.89\ \text{eV}。\]
换成波长（1 eV 对应约 \ud{1240}{nm}）：$1240 \div 1.89 \approx 656$ nm——一条红线。这正是氢光谱里最醒目的那条红线（叫 H$\alpha$ 线），天文望远镜靠它辨认宇宙里的氢云。一个写在纸上的公式，预言了恒星的颜色。有兴趣的话，把 $n=4$ 到 $n=2$、$n=5$ 到 $n=2$ 也算一算，对照氢光谱图看看它们落在什么颜色。跳过这一段不影响主线阅读。
\end{tiaozhan}

\section*{练一练}

\begin{sikaoti}
\item 画一幅氢原子的能级图（参照本章的阶梯图，$n=1$ 到 $n=4$），用向下箭头画出电子从 $n=3$ 跳回 $n=1$、以及从 $n=3$ 跳到 $n=2$ 再跳到 $n=1$ 的两条路径。哪一步吐出的光子能量最大？对应的光偏蓝还是偏红？图旁注明：横线代表能量值，不是电子在空间里的位置。
\item 用“楼层+户型+人数”的写法，写出氧（8个电子）和镁（12个电子）的电子排布，并回答：它们各自最外面一层住了几个电子？
\item 一个同学画电子云时，把 1s 电子画成绕核的一圈虚线，说“这是电子跑过的轨迹”。请用本章内容指出他的图错在哪两处，并把图改对（提示：改成一张“点越多表示越容易找到电子”的散布图）。
\item 烟花厂想做“紫色”烟花，师傅把锶盐（红）和铜盐（蓝）混在一起烧。从能级的角度解释这个配方为什么可能有效；再回答：能不能通过“把钠盐烧得更热”让钠的黄光变成蓝光？为什么？
\item 氦先在太阳上被发现、27年后才在地球上找到。这个故事里，科学家凭什么敢断言“这条线属于一种没人见过的元素”？如果当时有人反驳说“那只是仪器误差”，可以用什么实验来裁决？
\end{sikaoti}
